\section*{Цель}
Получение практических навыков по автоматизированному тестированию с помощью Selenium IDE.

\section*{Задачи}
\begin{enumerate}
    \item Установить браузерное расширение Selenium IDE.
    \item Выполнить упражнения для знакомства с возможностями Selenium IDE при тестировании клиентской части веб-приложений.
    \item Провести тестирование интерфейса самостоятельно выбранного веб-приложения.
    \item Зафиксировать результат обучения в отчете.
\end{enumerate}

\begin{figure}[H]
    \centering
    \includegraphics[width=1\textwidth]{figs/install.png}
    \caption{Подтверждение установленного Selenium IDE}
    \label{fig:install.png}
\end{figure}

\section{Примеры выполнения упражнений из раздела 4}

В данном разделе представлены результаты выполнения базовых упражнений по ознакомлению со средой Selenium IDE. Рассмотрены процессы записи тестовых сценариев для статической веб-страницы, модификации параметров поиска веб-элементов, проведения негативного тестирования с преднамеренно ошибочными значениями, а также сохранения и последующей загрузки программного кода автотеста.

\subsection{Разработка сценария тестирования на учебном веб-приложении}
\begin{figure}[H]
    \centering
    \includegraphics[width=\textwidth]{figs/4.1.1.png}
    \caption{Формы Command: click; Target: id=birch}
    \label{fig:4.1.1.png}
\end{figure}
\begin{figure}[H]
    \centering
    \includegraphics[width=\textwidth]{figs/4.1.2.png}
    \caption{Успешное завершение теста}
    \label{fig:4.1.2.png}
\end{figure}
\begin{figure}[H]
    \centering
    \includegraphics[width=\textwidth]{figs/4.1.3.png}
    \caption{Полное завершение теста}
    \label{fig:4.1.3.png}
\end{figure}

\begin{codemultipage}
    \captionof{listing}{Программный код строки с командой click().\label{lst:side1}}
    \nopagebreak
    \inputminted{side}{scripts/SelectorsTraining.side}
\end{codemultipage}

\subsection{Изменение параметров тестирования}
\begin{figure}[H]
    \centering
    \includegraphics[width=\textwidth]{figs/4.2.png}
    \caption{Изменение параметров поиска веб-элемента}
    \label{fig:4.2.png}
\end{figure}

\subsection{Негативное тестирование}
\begin{figure}[H]
    \centering
    \includegraphics[width=\textwidth]{figs/4.3.png}
    \caption{Результат выполнения негативного тестирования (тест провален)}
    \label{fig:4.3.png}
\end{figure}

\subsection{Открытие программного кода теста в Selenium IDE}
\begin{figure}[H]
    \centering
    \includegraphics[width=\textwidth]{figs/4.4.png}
    \caption{Загрузка ранее сохраненного проекта в среду Selenium IDE}
    \label{fig:4.4.png}
\end{figure}


\section{Предметная область выбранного для тестирования веб-приложения}
Для выполнения работы был выбран сайт компании «КЗ Групп», доступный по адресу \url{https://kzgroup.ru/}.

\textit{Предметная область.} Поставка промышленного оборудования и комплектующих для различных предприятий.

\textit{Назначение веб-приложения.} Сайт является визитной карточкой компании и содержит каталог продукции. Он предназначен для того, чтобы пользователи могли:
\begin{itemize}
    \item Посмотреть список доступного оборудования и его характеристики.
    \item Узнать о предоставляемых услугах.
    \item Найти контакты организации или оставить заявку на обратный звонок.
\end{itemize}


\section{Cценарий тестирования в табличной форме}
Для проверки работоспособности веб-приложения применяются следующие техники тестирования: дымовое тестирование, тестирование критического пути, а также метод «Черного ящика». 

Разработанный тестовый сценарий имитирует действия потенциального клиента, ищущего определенный тип оборудования.

\begin{longtblr}[
    caption = {Сценарий тестирования поиска и просмотра каталога},
    label = {tab:test_scenario},
]{
    colspec = { Q[c,m] X[l,m] X[l,m] },
    hlines, vlines,
    rowhead = 1,
}
    № & Последовательность действий & Ожидаемый результат \\
    1 & Открыть приложение в браузере (\url{https://kzgroup.ru/}) & Приложение открыто, загружена главная страница. \\
    2 & Проверить заголовок страницы & Заголовок содержит название компании. \\
    3 & Кликнуть по гамбургер меню & Гамбургер меню откроется. \\
    4 & Закрыть гамбургер меню & Гамбургер меню закроется. \\
    5 & Вернуться на главную страницу кликом по логотипу & Осуществлен возврат на стартовую страницу. \\
    6 & Перепроверить заголовок страницы & Заголовок содержит название компании. \\
    7 & Закрыть окно браузера & Окно закрыто. \\
\end{longtblr}


\section{Процесс тестирования}

% ИНСТРУКЦИЯ ПО ВЫПОЛНЕНИЮ ШАГОВ 6-8:
% 1. Открой Selenium IDE, создай новый проект "KZGroupTest".
% 2. Укажи Base URL: https://kzgroup.ru/
% 3. Нажми кнопку записи (Rec).
% 4. Выполни действия из таблицы выше в браузере:
%    - Дождись загрузки главной.
%    - Правый клик в пустом месте -> Selenium IDE -> Verify -> Title.
%    - Кликни на "Каталог".
%    - Найди поле поиска, введи "Насос", нажми Enter.
%    - Кликни на первый попавшийся насос.
%    - Правый клик по кнопке запроса цены/заказа -> Selenium IDE -> Assert -> Element Present (или Text).
%    - Кликни на логотип компании в шапке сайта.
% 5. Останови запись. Удали лишние шаги (например, случайные mouseOver), оставь только суть.
% 6. СДЕЛАЙ СКРИНШОТ №1: Окно Selenium IDE с записанным, но еще не запущенным кодом. Сохрани как figs/custom_test_raw.png.
% 7. Нажми кнопку Run current test. Дождись завершения.
% 8. Если тест упал (красный цвет) — выбери другой локатор (Target) в выпадающем списке для упавшего шага. Запусти снова.
% 9. СДЕЛАЙ СКРИНШОТ №2: Окно Selenium IDE после успешного прогона (все шаги зеленые, внизу виден лог с надписью completed successfully). Сохрани как figs/custom_test_done.png.
% 10. Сохрани проект (иконка дискеты) как kzgroup.side.

Процесс автоматизированного тестирования был зафиксирован в среде Selenium IDE. На \cref{fig:custom_test_raw.png} представлен программный код автотеста, полученный с помощью макрорекордера после выполнения действий пользователя согласно сценарию.

\begin{figure}[H]
    \centering
    \includegraphics[width=\textwidth]{figs/custom_test_raw.png}
    \caption{Программный код в среде Selenium IDE (до отладки)}
    \label{fig:custom_test_raw.png}
\end{figure}

На \cref{fig:custom_test_done.png} показан результат успешного выполнения отлаженного автотеста. В журнале логирования зафиксировано корректное прохождение всех проверок (verify/assert) и успешное завершение сценария.

\begin{figure}[H]
    \centering
    \includegraphics[width=\textwidth]{figs/custom_test_done.png}
    \caption{Отлаженный код автотеста и журнал логирования успешного выполнения}
    \label{fig:custom_test_done.png}
\end{figure}

% ИНСТРУКЦИЯ ПО ВЫПОЛНЕНИЮ ШАГА 10 (ПРИЛОЖЕНИЕ):
% Открой сохраненный файл kzgroup.side в любом текстовом редакторе (Блокнот, VS Code).
% Скопируй оттуда JSON-код твоего теста.
% Сохрани его в папку scripts/ под именем kzgroup.side.

\begin{codemultipage}
    \captionof{listing}{Фрагмент программного кода автотеста\label{lst:labkzgroup_sideel}}
    \nopagebreak
    \inputminted{json}{scripts/kzgroup.side}
\end{codemultipage}


\section{Выводы}
В ходе выполнения практической работы были получены навыки автоматизированного тестирования клиентской части веб-приложений с использованием инструмента Selenium IDE. Был успешно установлен и настроен соответствующий плагин для браузера, а также изучен базовый набор команд.

Основной задачей стала разработка собственного тестового сценария для сайта «КЗ Групп».

Практическая значимость работы заключается в освоении полного цикла создания автотеста: от проектирования сценария до записи, отладки локаторов и анализа журнала логирования. Полученные навыки позволяют автоматизировать рутинные проверки пользовательских интерфейсов и применять базовые методы регрессионного тестирования.